% Options for packages loaded elsewhere
\PassOptionsToPackage{unicode}{hyperref}
\PassOptionsToPackage{hyphens}{url}
%
\documentclass[
]{article}
\usepackage{amsmath,amssymb}
\usepackage{iftex}
\ifPDFTeX
  \usepackage[T1]{fontenc}
  \usepackage[utf8]{inputenc}
  \usepackage{textcomp} % provide euro and other symbols
\else % if luatex or xetex
  \usepackage{unicode-math} % this also loads fontspec
  \defaultfontfeatures{Scale=MatchLowercase}
  \defaultfontfeatures[\rmfamily]{Ligatures=TeX,Scale=1}
\fi
\usepackage{lmodern}
\ifPDFTeX\else
  % xetex/luatex font selection
\fi
% Use upquote if available, for straight quotes in verbatim environments
\IfFileExists{upquote.sty}{\usepackage{upquote}}{}
\IfFileExists{microtype.sty}{% use microtype if available
  \usepackage[]{microtype}
  \UseMicrotypeSet[protrusion]{basicmath} % disable protrusion for tt fonts
}{}
\makeatletter
\@ifundefined{KOMAClassName}{% if non-KOMA class
  \IfFileExists{parskip.sty}{%
    \usepackage{parskip}
  }{% else
    \setlength{\parindent}{0pt}
    \setlength{\parskip}{6pt plus 2pt minus 1pt}}
}{% if KOMA class
  \KOMAoptions{parskip=half}}
\makeatother
\usepackage{xcolor}
\usepackage[margin=1in]{geometry}
\usepackage{color}
\usepackage{fancyvrb}
\newcommand{\VerbBar}{|}
\newcommand{\VERB}{\Verb[commandchars=\\\{\}]}
\DefineVerbatimEnvironment{Highlighting}{Verbatim}{commandchars=\\\{\}}
% Add ',fontsize=\small' for more characters per line
\usepackage{framed}
\definecolor{shadecolor}{RGB}{248,248,248}
\newenvironment{Shaded}{\begin{snugshade}}{\end{snugshade}}
\newcommand{\AlertTok}[1]{\textcolor[rgb]{0.94,0.16,0.16}{#1}}
\newcommand{\AnnotationTok}[1]{\textcolor[rgb]{0.56,0.35,0.01}{\textbf{\textit{#1}}}}
\newcommand{\AttributeTok}[1]{\textcolor[rgb]{0.13,0.29,0.53}{#1}}
\newcommand{\BaseNTok}[1]{\textcolor[rgb]{0.00,0.00,0.81}{#1}}
\newcommand{\BuiltInTok}[1]{#1}
\newcommand{\CharTok}[1]{\textcolor[rgb]{0.31,0.60,0.02}{#1}}
\newcommand{\CommentTok}[1]{\textcolor[rgb]{0.56,0.35,0.01}{\textit{#1}}}
\newcommand{\CommentVarTok}[1]{\textcolor[rgb]{0.56,0.35,0.01}{\textbf{\textit{#1}}}}
\newcommand{\ConstantTok}[1]{\textcolor[rgb]{0.56,0.35,0.01}{#1}}
\newcommand{\ControlFlowTok}[1]{\textcolor[rgb]{0.13,0.29,0.53}{\textbf{#1}}}
\newcommand{\DataTypeTok}[1]{\textcolor[rgb]{0.13,0.29,0.53}{#1}}
\newcommand{\DecValTok}[1]{\textcolor[rgb]{0.00,0.00,0.81}{#1}}
\newcommand{\DocumentationTok}[1]{\textcolor[rgb]{0.56,0.35,0.01}{\textbf{\textit{#1}}}}
\newcommand{\ErrorTok}[1]{\textcolor[rgb]{0.64,0.00,0.00}{\textbf{#1}}}
\newcommand{\ExtensionTok}[1]{#1}
\newcommand{\FloatTok}[1]{\textcolor[rgb]{0.00,0.00,0.81}{#1}}
\newcommand{\FunctionTok}[1]{\textcolor[rgb]{0.13,0.29,0.53}{\textbf{#1}}}
\newcommand{\ImportTok}[1]{#1}
\newcommand{\InformationTok}[1]{\textcolor[rgb]{0.56,0.35,0.01}{\textbf{\textit{#1}}}}
\newcommand{\KeywordTok}[1]{\textcolor[rgb]{0.13,0.29,0.53}{\textbf{#1}}}
\newcommand{\NormalTok}[1]{#1}
\newcommand{\OperatorTok}[1]{\textcolor[rgb]{0.81,0.36,0.00}{\textbf{#1}}}
\newcommand{\OtherTok}[1]{\textcolor[rgb]{0.56,0.35,0.01}{#1}}
\newcommand{\PreprocessorTok}[1]{\textcolor[rgb]{0.56,0.35,0.01}{\textit{#1}}}
\newcommand{\RegionMarkerTok}[1]{#1}
\newcommand{\SpecialCharTok}[1]{\textcolor[rgb]{0.81,0.36,0.00}{\textbf{#1}}}
\newcommand{\SpecialStringTok}[1]{\textcolor[rgb]{0.31,0.60,0.02}{#1}}
\newcommand{\StringTok}[1]{\textcolor[rgb]{0.31,0.60,0.02}{#1}}
\newcommand{\VariableTok}[1]{\textcolor[rgb]{0.00,0.00,0.00}{#1}}
\newcommand{\VerbatimStringTok}[1]{\textcolor[rgb]{0.31,0.60,0.02}{#1}}
\newcommand{\WarningTok}[1]{\textcolor[rgb]{0.56,0.35,0.01}{\textbf{\textit{#1}}}}
\usepackage{graphicx}
\makeatletter
\newsavebox\pandoc@box
\newcommand*\pandocbounded[1]{% scales image to fit in text height/width
  \sbox\pandoc@box{#1}%
  \Gscale@div\@tempa{\textheight}{\dimexpr\ht\pandoc@box+\dp\pandoc@box\relax}%
  \Gscale@div\@tempb{\linewidth}{\wd\pandoc@box}%
  \ifdim\@tempb\p@<\@tempa\p@\let\@tempa\@tempb\fi% select the smaller of both
  \ifdim\@tempa\p@<\p@\scalebox{\@tempa}{\usebox\pandoc@box}%
  \else\usebox{\pandoc@box}%
  \fi%
}
% Set default figure placement to htbp
\def\fps@figure{htbp}
\makeatother
\setlength{\emergencystretch}{3em} % prevent overfull lines
\providecommand{\tightlist}{%
  \setlength{\itemsep}{0pt}\setlength{\parskip}{0pt}}
\setcounter{secnumdepth}{-\maxdimen} % remove section numbering
\usepackage{multirow}
\usepackage{multicol}
\usepackage{colortbl}
\usepackage[normalem]{ulem}
\usepackage{hhline}
\newlength\Oldarrayrulewidth
\newlength\Oldtabcolsep
\usepackage{longtable}
\usepackage{array}
\usepackage{hyperref}
\usepackage{float}
\usepackage{wrapfig}
\usepackage{bookmark}
\IfFileExists{xurl.sty}{\usepackage{xurl}}{} % add URL line breaks if available
\urlstyle{same}
\hypersetup{
  pdftitle={Mock Report: Northeast EPI},
  pdfauthor={Sam Long},
  hidelinks,
  pdfcreator={LaTeX via pandoc}}

\title{Mock Report: Northeast EPI}
\author{Sam Long}
\date{2026-09-25}

\begin{document}
\maketitle

\section{Introduction}\label{introduction}

This is an R markdown document. In R Studio, you can also generate
another kind of markdown file called an R notebook. The notebook allows
you to preview your work as you go in the viewer. This can also be
achieved, however, by checking ``Knit on Save'' when building a standard
markdown. Both approaches are excellent for report development and for
learning.

In this workshop, we'll cover the basics of R markdown and how to use it
to build automated reports. We'll cover:

\begin{itemize}
\item
  R markdown file format
\item
  Code chunk settings

  \begin{itemize}
  \item
    Output formatting
  \item
    Figure display
  \item
    Conditional execution
  \end{itemize}
\item
  Narrative options

  \begin{itemize}
  \item
    Embedded links
  \item
    Dynamic text
  \item
    Custom styling with CSS
  \end{itemize}
\item
  Parameters

  \begin{itemize}
  \item
    Automated report running
  \item
    Rendering via a script
  \item
    Error handling frameworks
  \end{itemize}
\item
  Useful packages for R markdown reports
\end{itemize}

Before we move on, take a look to the right side of the source pane and
you'll notice the Outline in a separate scrollable box. The Outline
allows you to jump to named sections and code chunks in your file.
Sections are created with ``\#'', and code chunks are named following
``\,```\{r'' in the chunk header. If you don't see the Outline, click
Tools \textgreater{} Global Options \textgreater{} R Markdown
\textgreater{} Basic and check ``Show document outline by default''.

\section{R Markdown Format}\label{r-markdown-format}

R Markdown files are specially formatted files that can be easily
``rendered'' to HTML files. This makes them a flexible option for
quickly converting code and notes into reports. Unlike R scripts, R
markdown files default to text. Unless the user creates a ``code
chunk'', anything written in the file (outside of YAML header) will be
interpreted as simple text. Several text styling options are available
in the default R markdown file; we'll cover these later. This section
and the Introduction above are narrative.

Data analyses and visualizations are incorporated via code chunks, which
are segments of the document specially formatted to interpret written
text as R code. Code chunks have adjustable settings that can be
globally and individually set for your document. We'll get into code
chunks below.

\section{Code Chunk Settings}\label{code-chunk-settings}

\begin{Shaded}
\begin{Highlighting}[]
\CommentTok{\# setting chunk options for my entire document }
\NormalTok{knitr}\SpecialCharTok{::}\NormalTok{opts\_chunk}\SpecialCharTok{$}\FunctionTok{set}\NormalTok{(}\AttributeTok{echo =} \ConstantTok{TRUE}\NormalTok{)}
\end{Highlighting}
\end{Shaded}

\begin{Shaded}
\begin{Highlighting}[]
\CommentTok{\# To ensure the document can be "Knitted" for a preview, packages must be}
\CommentTok{\# loaded in the document itself. If a document is rendered by a call in a }
\CommentTok{\# separate script, packages can be loaded elsewhere }
\ControlFlowTok{if}\NormalTok{ (}\SpecialCharTok{!}\FunctionTok{requireNamespace}\NormalTok{(}\StringTok{"flextable"}\NormalTok{, }\AttributeTok{quietly =} \ConstantTok{TRUE}\NormalTok{)) \{}
  \FunctionTok{install.packages}\NormalTok{(}\StringTok{"flextable"}\NormalTok{, }\AttributeTok{repos =} \StringTok{"https://r{-}project.org"}\NormalTok{)}
\NormalTok{\}}
\FunctionTok{library}\NormalTok{(flextable)}
\end{Highlighting}
\end{Shaded}

\begin{verbatim}
## Warning: package 'flextable' was built under R version 4.5.3
\end{verbatim}

\begin{Shaded}
\begin{Highlighting}[]
\FunctionTok{library}\NormalTok{(dplyr)}
\end{Highlighting}
\end{Shaded}

\begin{verbatim}
## Warning: package 'dplyr' was built under R version 4.5.3
\end{verbatim}

\begin{verbatim}
## 
## Attaching package: 'dplyr'
\end{verbatim}

\begin{verbatim}
## The following objects are masked from 'package:stats':
## 
##     filter, lag
\end{verbatim}

\begin{verbatim}
## The following objects are masked from 'package:base':
## 
##     intersect, setdiff, setequal, union
\end{verbatim}

\begin{Shaded}
\begin{Highlighting}[]
\FunctionTok{library}\NormalTok{(ggplot2)}
\end{Highlighting}
\end{Shaded}

\begin{verbatim}
## Warning: package 'ggplot2' was built under R version 4.5.3
\end{verbatim}

\begin{Shaded}
\begin{Highlighting}[]
\FunctionTok{library}\NormalTok{(readxl)}
\end{Highlighting}
\end{Shaded}

\begin{verbatim}
## Warning: package 'readxl' was built under R version 4.5.3
\end{verbatim}

\begin{Shaded}
\begin{Highlighting}[]
\FunctionTok{library}\NormalTok{(flextable)}
\end{Highlighting}
\end{Shaded}

\subsection{Chunk Arguments for Text
Output}\label{chunk-arguments-for-text-output}

Here we'll demonstrate a few of the code chunk options available in
markdown:

\subsubsection{Message}\label{message}

The \texttt{message} argument sets whether to display message outputs in
rendered file from functions executed in code. Messages are generated by
functions to alert users to certain processes.

\begin{Shaded}
\begin{Highlighting}[]
\CommentTok{\# Custom function that sends a message }
\NormalTok{count }\OtherTok{\textless{}{-}} \FunctionTok{c}\NormalTok{(}\DecValTok{1}\NormalTok{,}\DecValTok{2}\NormalTok{,}\DecValTok{3}\NormalTok{)}
\NormalTok{send\_message }\OtherTok{\textless{}{-}} \ControlFlowTok{function}\NormalTok{(vector) \{}
  \ControlFlowTok{for}\NormalTok{ (i }\ControlFlowTok{in} \FunctionTok{seq\_along}\NormalTok{(vector)) \{}
    \CommentTok{\# count down!}
    \FunctionTok{message}\NormalTok{(i, }\StringTok{" Mississippi!"}\NormalTok{)}
\NormalTok{  \}}
\NormalTok{\}}
\end{Highlighting}
\end{Shaded}

\begin{Shaded}
\begin{Highlighting}[]
\CommentTok{\# message prints out separately }
\FunctionTok{send\_message}\NormalTok{(count)}
\end{Highlighting}
\end{Shaded}

\begin{verbatim}
## 1 Mississippi!
\end{verbatim}

\begin{verbatim}
## 2 Mississippi!
\end{verbatim}

\begin{verbatim}
## 3 Mississippi!
\end{verbatim}

\begin{Shaded}
\begin{Highlighting}[]
\CommentTok{\# keeps message from this function hidden}
\FunctionTok{send\_message}\NormalTok{(count)}
\end{Highlighting}
\end{Shaded}

\subsubsection{Error \& Warning}\label{error-warning}

The \texttt{error} argument sets whether to continue rendering if an
error is encountered (FALSE by default), and the \texttt{warning}
argument sets whether to display warning outputs in rendered file from
functions executed in code (TRUE by default).

\begin{Shaded}
\begin{Highlighting}[]
\CommentTok{\# this will throw an error that prevents complete rendering of the file}
\CommentTok{\# unless error = T}
\FunctionTok{sum}\NormalTok{(}\FunctionTok{c}\NormalTok{(}\DecValTok{1}\NormalTok{, }\DecValTok{2}\NormalTok{, }\StringTok{"3"}\NormalTok{))}
\end{Highlighting}
\end{Shaded}

\begin{verbatim}
## Error in sum(c(1, 2, "3")): invalid 'type' (character) of argument
\end{verbatim}

\begin{Shaded}
\begin{Highlighting}[]
\CommentTok{\# this will only print if the error is "ignored"}
\FunctionTok{print}\NormalTok{(}\StringTok{"Nothing to see here, folks!"}\NormalTok{)}
\end{Highlighting}
\end{Shaded}

\begin{verbatim}
## [1] "Nothing to see here, folks!"
\end{verbatim}

\begin{Shaded}
\begin{Highlighting}[]
\CommentTok{\# this function produces a warning message that NAs were imputed}
\CommentTok{\# for all values where coercion to numeric was not possible}
\FunctionTok{as.numeric}\NormalTok{(}\FunctionTok{c}\NormalTok{(}\StringTok{"1"}\NormalTok{, }\StringTok{"2"}\NormalTok{, }\StringTok{"three"}\NormalTok{))}
\end{Highlighting}
\end{Shaded}

\begin{verbatim}
## Warning: NAs introduced by coercion
\end{verbatim}

\begin{verbatim}
## [1]  1  2 NA
\end{verbatim}

\begin{Shaded}
\begin{Highlighting}[]
\CommentTok{\# warning = F will hide the warning in the output}
\FunctionTok{as.numeric}\NormalTok{(}\FunctionTok{c}\NormalTok{(}\StringTok{"1"}\NormalTok{, }\StringTok{"2"}\NormalTok{, }\StringTok{"three"}\NormalTok{))}
\end{Highlighting}
\end{Shaded}

\begin{verbatim}
## [1]  1  2 NA
\end{verbatim}

\subsubsection{Include}\label{include}

The \texttt{include} argument sets whether to display chunk results and
the chunk itself in output document (TRUE by default). So far, we've
chosen to display all our text outputs, but this argument makes it easy
to obscure them.

\subsubsection{Results}\label{results}

The \texttt{results} argument has two settings, ``markup'', where text
is displayed in boxes (see all the outputs above), and ``asis'' where
text output is converted to plain text.

In combination with the cat() function, this argument allows for
narrative/HTML sections to be dynamically included in the output.

\begin{Shaded}
\begin{Highlighting}[]
\FunctionTok{print}\NormalTok{(}\StringTok{"help, I\textquotesingle{}m stuck in this box!"}\NormalTok{)}
\end{Highlighting}
\end{Shaded}

\begin{verbatim}
## [1] "help, I'm stuck in this box!"
\end{verbatim}

\begin{Shaded}
\begin{Highlighting}[]
\CommentTok{\# present text output as a narrative line}
\FunctionTok{print}\NormalTok{(}\StringTok{"Ah, much better."}\NormalTok{)}
\end{Highlighting}
\end{Shaded}

{[}1{]} ``Ah, much better.''

\begin{Shaded}
\begin{Highlighting}[]
\CommentTok{\# use cat() for even cleaner text}
\FunctionTok{cat}\NormalTok{(}\StringTok{"Now I can do anything!!"}\NormalTok{)}
\end{Highlighting}
\end{Shaded}

Now I can do anything!!

\subsection{Chunk Evaluation}\label{chunk-evaluation}

You can choose whether or not to actually execute a code chunk by
setting the \texttt{eval} argument, which takes a Boolean operator. Hard
coding this operator is useful during development, but setting it
dynamically can really come in handy when rendering automated reports.

\begin{Shaded}
\begin{Highlighting}[]
\CommentTok{\# say we only want a certain chunk to evaluate if a certain user}
\CommentTok{\# is running it}

\CommentTok{\# Hey user, what\textquotesingle{}s your name??}
\NormalTok{my\_name }\OtherTok{\textless{}{-}} \StringTok{"Sam"}

\CommentTok{\# set the value }
\ControlFlowTok{if}\NormalTok{ (my\_name }\SpecialCharTok{==} \StringTok{"Sam"}\NormalTok{) \{}
\NormalTok{  name\_val }\OtherTok{\textless{}{-}}\NormalTok{ T}
\NormalTok{\} }\ControlFlowTok{else}\NormalTok{ \{}
\NormalTok{  name\_val }\OtherTok{\textless{}{-}}\NormalTok{ F}
\NormalTok{\}}
\end{Highlighting}
\end{Shaded}

Hey Sam, don't forget to save the report to the correct folder!

\subsection{Chunk Arguments for Figures and
Plots}\label{chunk-arguments-for-figures-and-plots}

There are many ways to format figures and plots in reports. It can be
very convenient to apply certain formatting choices directly in the
chunk settings, though more complex formatting must be achieved through
other means.

\subsubsection{Embedding Figures}\label{embedding-figures}

Pre-saved images can be included in outputs via calls in code chunks.
You can also use chunk arguments to provide captions and labels.

\begin{Shaded}
\begin{Highlighting}[]
\CommentTok{\# if the image is not in the working directory, be sure to provide the full path}
\NormalTok{knitr}\SpecialCharTok{::}\FunctionTok{include\_graphics}\NormalTok{(}\StringTok{\textquotesingle{}ok\_cat.jpg\textquotesingle{}}\NormalTok{)}
\end{Highlighting}
\end{Shaded}

\begin{figure}
\centering
\pandocbounded{\includegraphics[keepaspectratio]{ok_cat.jpg}}
\caption{The cat is just doing his best}
\end{figure}

\subsubsection{Resizing and Placing
Figures}\label{resizing-and-placing-figures}

It's also possible to control image size and alignment via chunk
arguments.You can call \texttt{out.height} and \texttt{out.width} to
specify the size of the image. These arguments work with saved images,
and will not be affected by the similarly named \texttt{fig.height},
\texttt{fig.width}, and \texttt{fig.dim}, which work only when calling R
objects (such as a ggplot). These arguments interpret the value input in
inches.

The \texttt{fig.align} argument allows you to quickly re-position the
image, which can be helpful for aesthetics of your report. There are
more ways to adjust figure size and position, and we'll address some of
these later on.

\begin{Shaded}
\begin{Highlighting}[]
\CommentTok{\# when using out.height and out.width, you can use either the percent of }
\CommentTok{\# available space you\textquotesingle{}d like the image to take up or a specific measurement}
\CommentTok{\# such as pixels}
\NormalTok{knitr}\SpecialCharTok{::}\FunctionTok{include\_graphics}\NormalTok{(}\StringTok{\textquotesingle{}ok\_cat.jpg\textquotesingle{}}\NormalTok{)}
\end{Highlighting}
\end{Shaded}

\begin{figure}
\includegraphics[width=700px,height=20px]{ok_cat} \caption{Now the cat is flat}\label{fig:show_fig_wide}
\end{figure}

\begin{Shaded}
\begin{Highlighting}[]
\CommentTok{\# fig.align works with saved images and R objects alike!}
\NormalTok{knitr}\SpecialCharTok{::}\FunctionTok{include\_graphics}\NormalTok{(}\StringTok{\textquotesingle{}ok\_cat.jpg\textquotesingle{}}\NormalTok{)}
\end{Highlighting}
\end{Shaded}

\begin{figure}

\hfill{}\includegraphics[width=200px]{ok_cat} 

\caption{Now the cat is right}\label{fig:show_fig_right}
\end{figure}

\section{Narrative Options}\label{narrative-options}

\subsection{Type Face Adjustments and
Paragraphs}\label{type-face-adjustments-and-paragraphs}

Markdown files understand certain cues in normal text that will change
the face of the text. Here we'll also bring attention the easily applied
paragraph delineations, which have been used throughout this document
thus far. Pressing enter will send your text to a new line, but pressing
it twice will start a new paragraph, as shown below.

Bold text can be achieved with \textbf{double asterisks}, or
alternatively with \textbf{double underscores}.

Italics can be applied with \emph{single asterisks} and \emph{single
underscores}.

You can also simultaneously apply both with \textbf{\emph{triple
asterisks}} or \textbf{\emph{underscores}}.

Generally speaking, asterisks behave more stably than underlines, and
can even be used to just change part of a word, like t\textbf{hi}s,
whereas the underscores don't work in t\_\_hi\_\_s case.

It is also possible to force markdown to interpret these characters
simply, and not as type face adjustments. To do this, we need to use an
``escape character''. In R markdown, the escape character is the
backslash. Here, I can ``escape'' the effect of my *asterisks* to keep
the font face normal and show the asterisks themselves.

\subsection{Embedding Links, Dynamic Text, and
Figures}\label{embedding-links-dynamic-text-and-figures}

Adding links to your text is very easy. Simply type the visible text
within square brackets immediately followed by the link in parentheses.
If I want to direct viewers to a brief overview of the cat meme I used
above, I'd link to it
\href{https://en.meming.world/wiki/Thumbs_Up_Crying_Cat}{here}.

It's also common to want to include text that changes based on input
parameters or the outcome of functions called in code chunks. You can do
this by calling objects ``inline'', where markdown will interpret them
as if they are in code chunks. Earlier, I set the object
\texttt{my\_name}, I can interpret it dynamically by using back ticks
and a lower case r:

``Hi, my name is Sam. Nice to meet you!''

You can also write new code to be evaluated in these inline chunks,
though it's best to keep these calls brief. A great example is using the
inline to dynamically write today's date.

Today's date is November 06, 2026.

Finally, while it's preferable to add figures in other ways, it is
possible to embed figures in text by using the link format preceded by
an exclamation mark. Here's a different cat:
\pandocbounded{\includegraphics[keepaspectratio]{cat_table.jpg}}.

\subsection{Lists, Block Quotes, and Horizontal
Rules}\label{lists-block-quotes-and-horizontal-rules}

Adding bullets and lists to markdown is fairly simple:

\begin{enumerate}
\def\labelenumi{\arabic{enumi}.}
\tightlist
\item
  Press enter twice to start a new paragraph
\item
  Type ``1.'' to start your numbered list
\item
  Repeat as necessary

  \begin{enumerate}
  \def\labelenumii{\alph{enumii}.}
  \tightlist
  \item
    Simply indent \textbf{twice} and type sub-items as you'd like them
    to appear
  \end{enumerate}
\end{enumerate}

\begin{itemize}
\tightlist
\item
  Bullets can be added using dashes
\item
  You can also use ``+'' or ``*'' but it's best to pick a character and
  stick to it!
\end{itemize}

Block quotes can be added using the right arrow character,
``\textgreater{}''. Let's say I want to quote the article I linked
above. I'd do it by starting a new paragraph, indenting once, and using
the arrow:

\begin{quote}
``The ``Thumbs Up Crying Cat'' meme is a reaction image involving the
photoshopped image of a cat giving the viewer a thumbs-up gesture. This
sister image to ``Crying Cat'' has commonly been used to indicate a
mixture of sadness but acceptance of a situation, possibly to avoid
being difficult.''
\end{quote}

To end this section with a nice horizontal rule, I'll simply type three
lone underscores in a new paragraph.

\begin{center}\rule{0.5\linewidth}{0.5pt}\end{center}

\subsection{Custom Styling with CSS}\label{custom-styling-with-css}

For the most precise control over fonts, text appearance, and
arrangement of elements in a document, you can use a .css file to
``style'' your markdown. The CSS file will contain specially formatted
instructions which you can then call in your markdown file to apply.
Creating a new CSS file in R is as easy as clicking File \textgreater{}
New file \textgreater{} CSS File. Once created, save it with a name and
then link your markdown to the CSS in the YAML header.

In the following sections I'll show how CSS can be incorporated to
achieve great visual results. There are several ways to do this, but my
team likes to use ``div'' elements and grid layouts to create products.
When using ``div'' elements, you have more control, but that also means
fewer things happen automatically. Headings, for example need to be
``styled'' and then called with the ``class'' argument in the div
opener:

Here's My CSS Styled Heading

CSS is great for applying styling options to a whole section of text.
Rather than wrapping all of this in asterisks, I applied my custom CSS
style. Also, within div elements we do still need to use escape
characters to simply type **special characters**.

With CSS and the grid layout, we can design page layouts with
``containers'' and use these to arrange blocks of text or images on the
page. The container will be the outermost div element, and the objects
within the container will be in nested div elements. See how this looks
below.

``One cat just leads to another''

- Ernest Hemingway

Total customization via CSS is great, but specifying everything can be
very time-consuming. R markdown documents use Bootstrap styling by
default, and thus can have a number of pre-built ``Bootswatch themes''
applied to them. The \texttt{bslib} package gives you access to the
widest range of standard and modern
\href{https://bootswatch.com/}{Bootswatch themes}.

\begin{verbatim}
## Warning: package 'bookdown' was built under R version 4.5.3
\end{verbatim}

Finally, if the desired output format is simply HTML, then you do not
need to worry about pagination. However, if the desired output of your
document is PDF, pagination is key. My team uses the \texttt{pagedown}
package to handle paging and HTML to PDF conversion.

If you do not specify page breaks in your markdown file and convert the
rendered HTML to PDF, breaks will automatically insert as pages fill up.
This can be aesthetically pleasing by chance, but if the goal is to
create an automated report with predictable formatting and page content,
then setting page breaks is a \emph{must}. I do this below by typing
backslash pagebreak:

\pagebreak

Here's a custom header

Hello there Sam!

\section{Useful Packages for Markdown
Reports}\label{useful-packages-for-markdown-reports}

Here I'll use the \texttt{flextable} package to create a nice table of
useful packages for writing attractive HTML reports in markdown. This is
by no means an exhaustive list, but my team returns to these packages
time and time again when building data products.

\begin{verbatim}
## Warning: fonts used in `flextable` are ignored because the `pdflatex` engine is
## used and not `xelatex` or `lualatex`. You can avoid this warning by using
## `set_flextable_defaults(fonts_ignore=TRUE)` or use a compatible engine (e.g.
## `pdf-engine: xelatex`).
\end{verbatim}

\global\setlength{\Oldarrayrulewidth}{\arrayrulewidth}

\global\setlength{\Oldtabcolsep}{\tabcolsep}

\setlength{\tabcolsep}{2pt}

\renewcommand*{\arraystretch}{1.5}



\providecommand{\ascline}[3]{\noalign{\global\arrayrulewidth #1}\arrayrulecolor[HTML]{#2}\cline{#3}}

\begin{longtable}[c]{|p{1.50in}|p{1.50in}|p{1.50in}|p{1.50in}}



\hhline{>{\arrayrulecolor[HTML]{666666}\global\arrayrulewidth=1.5pt}->{\arrayrulecolor[HTML]{666666}\global\arrayrulewidth=1.5pt}->{\arrayrulecolor[HTML]{666666}\global\arrayrulewidth=1.5pt}->{\arrayrulecolor[HTML]{666666}\global\arrayrulewidth=1.5pt}-}

\multicolumn{4}{>{\raggedright}m{\dimexpr 6in+6\tabcolsep}}{\advance\leftskip by 5.0pt\relax \advance\rightskip by 5.0pt\relax \textcolor[HTML]{000000}{\fontsize{11}{11}\selectfont{\textbf{Useful\ Packages\ for\ Markdown\ Reports}}}} \\

\noalign{\global\arrayrulewidth 0pt}\arrayrulecolor[HTML]{000000}

\hhline{>{\arrayrulecolor[HTML]{666666}\global\arrayrulewidth=1.5pt}->{\arrayrulecolor[HTML]{666666}\global\arrayrulewidth=1.5pt}->{\arrayrulecolor[HTML]{666666}\global\arrayrulewidth=1.5pt}->{\arrayrulecolor[HTML]{666666}\global\arrayrulewidth=1.5pt}-}



\multicolumn{1}{>{\raggedright}m{\dimexpr 1.5in+0\tabcolsep}}{\advance\leftskip by 5.0pt\relax \advance\rightskip by 5.0pt\relax \textcolor[HTML]{000000}{\fontsize{11}{11}\selectfont{\textbf{Package(s)}}}} & \multicolumn{1}{>{\raggedright}m{\dimexpr 1.5in+0\tabcolsep}}{\advance\leftskip by 5.0pt\relax \advance\rightskip by 5.0pt\relax \textcolor[HTML]{000000}{\fontsize{11}{11}\selectfont{\textbf{Description}}}} & \multicolumn{1}{>{\raggedright}m{\dimexpr 1.5in+0\tabcolsep}}{\advance\leftskip by 5.0pt\relax \advance\rightskip by 5.0pt\relax \textcolor[HTML]{000000}{\fontsize{11}{11}\selectfont{\textbf{Pros}}}} & \multicolumn{1}{>{\raggedright}m{\dimexpr 1.5in+0\tabcolsep}}{\advance\leftskip by 5.0pt\relax \advance\rightskip by 5.0pt\relax \textcolor[HTML]{000000}{\fontsize{11}{11}\selectfont{\textbf{Cons}}}} \\

\noalign{\global\arrayrulewidth 0pt}\arrayrulecolor[HTML]{000000}

\hhline{>{\arrayrulecolor[HTML]{666666}\global\arrayrulewidth=1.5pt}->{\arrayrulecolor[HTML]{666666}\global\arrayrulewidth=1.5pt}->{\arrayrulecolor[HTML]{666666}\global\arrayrulewidth=1.5pt}->{\arrayrulecolor[HTML]{666666}\global\arrayrulewidth=1.5pt}-}\endfirsthead 

\hhline{>{\arrayrulecolor[HTML]{666666}\global\arrayrulewidth=1.5pt}->{\arrayrulecolor[HTML]{666666}\global\arrayrulewidth=1.5pt}->{\arrayrulecolor[HTML]{666666}\global\arrayrulewidth=1.5pt}->{\arrayrulecolor[HTML]{666666}\global\arrayrulewidth=1.5pt}-}

\multicolumn{4}{>{\raggedright}m{\dimexpr 6in+6\tabcolsep}}{\advance\leftskip by 5.0pt\relax \advance\rightskip by 5.0pt\relax \textcolor[HTML]{000000}{\fontsize{11}{11}\selectfont{\textbf{Useful\ Packages\ for\ Markdown\ Reports}}}} \\

\noalign{\global\arrayrulewidth 0pt}\arrayrulecolor[HTML]{000000}

\hhline{>{\arrayrulecolor[HTML]{666666}\global\arrayrulewidth=1.5pt}->{\arrayrulecolor[HTML]{666666}\global\arrayrulewidth=1.5pt}->{\arrayrulecolor[HTML]{666666}\global\arrayrulewidth=1.5pt}->{\arrayrulecolor[HTML]{666666}\global\arrayrulewidth=1.5pt}-}



\multicolumn{1}{>{\raggedright}m{\dimexpr 1.5in+0\tabcolsep}}{\advance\leftskip by 5.0pt\relax \advance\rightskip by 5.0pt\relax \textcolor[HTML]{000000}{\fontsize{11}{11}\selectfont{\textbf{Package(s)}}}} & \multicolumn{1}{>{\raggedright}m{\dimexpr 1.5in+0\tabcolsep}}{\advance\leftskip by 5.0pt\relax \advance\rightskip by 5.0pt\relax \textcolor[HTML]{000000}{\fontsize{11}{11}\selectfont{\textbf{Description}}}} & \multicolumn{1}{>{\raggedright}m{\dimexpr 1.5in+0\tabcolsep}}{\advance\leftskip by 5.0pt\relax \advance\rightskip by 5.0pt\relax \textcolor[HTML]{000000}{\fontsize{11}{11}\selectfont{\textbf{Pros}}}} & \multicolumn{1}{>{\raggedright}m{\dimexpr 1.5in+0\tabcolsep}}{\advance\leftskip by 5.0pt\relax \advance\rightskip by 5.0pt\relax \textcolor[HTML]{000000}{\fontsize{11}{11}\selectfont{\textbf{Cons}}}} \\

\noalign{\global\arrayrulewidth 0pt}\arrayrulecolor[HTML]{000000}

\hhline{>{\arrayrulecolor[HTML]{666666}\global\arrayrulewidth=1.5pt}->{\arrayrulecolor[HTML]{666666}\global\arrayrulewidth=1.5pt}->{\arrayrulecolor[HTML]{666666}\global\arrayrulewidth=1.5pt}->{\arrayrulecolor[HTML]{666666}\global\arrayrulewidth=1.5pt}-}\endhead



\multicolumn{4}{>{\cellcolor[HTML]{9DD0FA}\raggedright}m{\dimexpr 6in+6\tabcolsep}}{\advance\leftskip by 5.0pt\relax \advance\rightskip by 5.0pt\relax \textcolor[HTML]{000000}{\fontsize{11}{11}\selectfont{Tables}}} \\

\noalign{\global\arrayrulewidth 0pt}\arrayrulecolor[HTML]{000000}





\multicolumn{1}{>{\cellcolor[HTML]{C7E4FC}\raggedright}m{\dimexpr 1.5in+0\tabcolsep}}{\advance\leftskip by 5.0pt\relax \advance\rightskip by 5.0pt\relax \textcolor[HTML]{000000}{\fontsize{11}{11}\selectfont{kable,\ kableExtra}}} & \multicolumn{1}{>{\cellcolor[HTML]{C7E4FC}\raggedright}m{\dimexpr 1.5in+0\tabcolsep}}{\advance\leftskip by 5.0pt\relax \advance\rightskip by 5.0pt\relax \textcolor[HTML]{000000}{\fontsize{11}{11}\selectfont{Highly\ customizable\ table\ formatting\ for\ HTML}}} & \multicolumn{1}{>{\cellcolor[HTML]{C7E4FC}\raggedright}m{\dimexpr 1.5in+0\tabcolsep}}{\advance\leftskip by 5.0pt\relax \advance\rightskip by 5.0pt\relax \textcolor[HTML]{000000}{\fontsize{11}{11}\selectfont{Intuitive\ syntax,\ grouped\ columns,\ mid-table\ spanners}}} & \multicolumn{1}{>{\cellcolor[HTML]{C7E4FC}\raggedright}m{\dimexpr 1.5in+0\tabcolsep}}{\advance\leftskip by 5.0pt\relax \advance\rightskip by 5.0pt\relax \textcolor[HTML]{000000}{\fontsize{11}{11}\selectfont{Doesn't\ play\ nice\ with\ pagedown}}} \\

\noalign{\global\arrayrulewidth 0pt}\arrayrulecolor[HTML]{000000}





\multicolumn{1}{>{\cellcolor[HTML]{C7E4FC}\raggedright}m{\dimexpr 1.5in+0\tabcolsep}}{\advance\leftskip by 5.0pt\relax \advance\rightskip by 5.0pt\relax \textcolor[HTML]{000000}{\fontsize{11}{11}\selectfont{flextable}}} & \multicolumn{1}{>{\cellcolor[HTML]{C7E4FC}\raggedright}m{\dimexpr 1.5in+0\tabcolsep}}{\advance\leftskip by 5.0pt\relax \advance\rightskip by 5.0pt\relax \textcolor[HTML]{000000}{\fontsize{11}{11}\selectfont{Highly\ cusotmizable\ table\ formatting\ for\ PDF/PowerPoint/Word}}} & \multicolumn{1}{>{\cellcolor[HTML]{C7E4FC}\raggedright}m{\dimexpr 1.5in+0\tabcolsep}}{\advance\leftskip by 5.0pt\relax \advance\rightskip by 5.0pt\relax \textcolor[HTML]{000000}{\fontsize{11}{11}\selectfont{The\ easiest\ package\ for\ formatting\ tables\ beyond\ HTML}}} & \multicolumn{1}{>{\cellcolor[HTML]{C7E4FC}\raggedright}m{\dimexpr 1.5in+0\tabcolsep}}{\advance\leftskip by 5.0pt\relax \advance\rightskip by 5.0pt\relax \textcolor[HTML]{000000}{\fontsize{11}{11}\selectfont{Less\ flexible\ and\ intuitive\ than\ kable}}} \\

\noalign{\global\arrayrulewidth 0pt}\arrayrulecolor[HTML]{000000}





\multicolumn{1}{>{\cellcolor[HTML]{C7E4FC}\raggedright}m{\dimexpr 1.5in+0\tabcolsep}}{\advance\leftskip by 5.0pt\relax \advance\rightskip by 5.0pt\relax \textcolor[HTML]{000000}{\fontsize{11}{11}\selectfont{DT}}} & \multicolumn{1}{>{\cellcolor[HTML]{C7E4FC}\raggedright}m{\dimexpr 1.5in+0\tabcolsep}}{\advance\leftskip by 5.0pt\relax \advance\rightskip by 5.0pt\relax \textcolor[HTML]{000000}{\fontsize{11}{11}\selectfont{Interactive\ tables\ for\ HTML}}} & \multicolumn{1}{>{\cellcolor[HTML]{C7E4FC}\raggedright}m{\dimexpr 1.5in+0\tabcolsep}}{\advance\leftskip by 5.0pt\relax \advance\rightskip by 5.0pt\relax \textcolor[HTML]{000000}{\fontsize{11}{11}\selectfont{Easily\ create\ filterable\ and\ sortable\ tables,\ even\ with\ large\ data}}} & \multicolumn{1}{>{\cellcolor[HTML]{C7E4FC}\raggedright}m{\dimexpr 1.5in+0\tabcolsep}}{\advance\leftskip by 5.0pt\relax \advance\rightskip by 5.0pt\relax \textcolor[HTML]{000000}{\fontsize{11}{11}\selectfont{Relatively\ few\ customization\ options}}} \\

\noalign{\global\arrayrulewidth 0pt}\arrayrulecolor[HTML]{000000}





\multicolumn{1}{>{\cellcolor[HTML]{C7E4FC}\raggedright}m{\dimexpr 1.5in+0\tabcolsep}}{\advance\leftskip by 5.0pt\relax \advance\rightskip by 5.0pt\relax \textcolor[HTML]{000000}{\fontsize{11}{11}\selectfont{reactable}}} & \multicolumn{1}{>{\cellcolor[HTML]{C7E4FC}\raggedright}m{\dimexpr 1.5in+0\tabcolsep}}{\advance\leftskip by 5.0pt\relax \advance\rightskip by 5.0pt\relax \textcolor[HTML]{000000}{\fontsize{11}{11}\selectfont{Interactive\ and\ customizable\ tables\ for\ HTML}}} & \multicolumn{1}{>{\cellcolor[HTML]{C7E4FC}\raggedright}m{\dimexpr 1.5in+0\tabcolsep}}{\advance\leftskip by 5.0pt\relax \advance\rightskip by 5.0pt\relax \textcolor[HTML]{000000}{\fontsize{11}{11}\selectfont{More\ formatting\ and\ interactive\ options\ than\ DT\ with\ great\ built-in\ templates}}} & \multicolumn{1}{>{\cellcolor[HTML]{C7E4FC}\raggedright}m{\dimexpr 1.5in+0\tabcolsep}}{\advance\leftskip by 5.0pt\relax \advance\rightskip by 5.0pt\relax \textcolor[HTML]{000000}{\fontsize{11}{11}\selectfont{Syntax\ for\ specific\ cusotmization\ is\ challenging}}} \\

\noalign{\global\arrayrulewidth 0pt}\arrayrulecolor[HTML]{000000}





\multicolumn{4}{>{\cellcolor[HTML]{89F781}\raggedright}m{\dimexpr 6in+6\tabcolsep}}{\advance\leftskip by 5.0pt\relax \advance\rightskip by 5.0pt\relax \textcolor[HTML]{000000}{\fontsize{11}{11}\selectfont{Visualizations}}} \\

\noalign{\global\arrayrulewidth 0pt}\arrayrulecolor[HTML]{000000}





\multicolumn{1}{>{\cellcolor[HTML]{C6FCC2}\raggedright}m{\dimexpr 1.5in+0\tabcolsep}}{\advance\leftskip by 5.0pt\relax \advance\rightskip by 5.0pt\relax \textcolor[HTML]{000000}{\fontsize{11}{11}\selectfont{ggplot2}}} & \multicolumn{1}{>{\cellcolor[HTML]{C6FCC2}\raggedright}m{\dimexpr 1.5in+0\tabcolsep}}{\advance\leftskip by 5.0pt\relax \advance\rightskip by 5.0pt\relax \textcolor[HTML]{000000}{\fontsize{11}{11}\selectfont{Highly\ customizable\ plots\ and\ figures}}} & \multicolumn{1}{>{\cellcolor[HTML]{C6FCC2}\raggedright}m{\dimexpr 1.5in+0\tabcolsep}}{\advance\leftskip by 5.0pt\relax \advance\rightskip by 5.0pt\relax \textcolor[HTML]{000000}{\fontsize{11}{11}\selectfont{Intuitive\ layering\ approach\ to\ building\ plots,\ plethora\ of\ add-on\ packages\ for\ expanded\ utility}}} & \multicolumn{1}{>{\cellcolor[HTML]{C6FCC2}\raggedright}m{\dimexpr 1.5in+0\tabcolsep}}{\advance\leftskip by 5.0pt\relax \advance\rightskip by 5.0pt\relax \textcolor[HTML]{000000}{\fontsize{11}{11}\selectfont{Final\ plot\ code\ can\ be\ verbose,\ conflicting\ arguments\ can\ be\ frustrating\ to\ resolve}}} \\

\noalign{\global\arrayrulewidth 0pt}\arrayrulecolor[HTML]{000000}





\multicolumn{1}{>{\cellcolor[HTML]{C6FCC2}\raggedright}m{\dimexpr 1.5in+0\tabcolsep}}{\advance\leftskip by 5.0pt\relax \advance\rightskip by 5.0pt\relax \textcolor[HTML]{000000}{\fontsize{11}{11}\selectfont{patchwork}}} & \multicolumn{1}{>{\cellcolor[HTML]{C6FCC2}\raggedright}m{\dimexpr 1.5in+0\tabcolsep}}{\advance\leftskip by 5.0pt\relax \advance\rightskip by 5.0pt\relax \textcolor[HTML]{000000}{\fontsize{11}{11}\selectfont{Arrangement\ of\ plots\ into\ composite\ figures}}} & \multicolumn{1}{>{\cellcolor[HTML]{C6FCC2}\raggedright}m{\dimexpr 1.5in+0\tabcolsep}}{\advance\leftskip by 5.0pt\relax \advance\rightskip by 5.0pt\relax \textcolor[HTML]{000000}{\fontsize{11}{11}\selectfont{Simple\ syntax\ for\ arranging\ plots\ and\ collapsing\ labels}}} & \multicolumn{1}{>{\cellcolor[HTML]{C6FCC2}\raggedright}m{\dimexpr 1.5in+0\tabcolsep}}{\advance\leftskip by 5.0pt\relax \advance\rightskip by 5.0pt\relax \textcolor[HTML]{000000}{\fontsize{11}{11}\selectfont{None}}} \\

\noalign{\global\arrayrulewidth 0pt}\arrayrulecolor[HTML]{000000}





\multicolumn{1}{>{\cellcolor[HTML]{C6FCC2}\raggedright}m{\dimexpr 1.5in+0\tabcolsep}}{\advance\leftskip by 5.0pt\relax \advance\rightskip by 5.0pt\relax \textcolor[HTML]{000000}{\fontsize{11}{11}\selectfont{sf}}} & \multicolumn{1}{>{\cellcolor[HTML]{C6FCC2}\raggedright}m{\dimexpr 1.5in+0\tabcolsep}}{\advance\leftskip by 5.0pt\relax \advance\rightskip by 5.0pt\relax \textcolor[HTML]{000000}{\fontsize{11}{11}\selectfont{Comprehensive\ spatial\ analysis\ and\ mapping}}} & \multicolumn{1}{>{\cellcolor[HTML]{C6FCC2}\raggedright}m{\dimexpr 1.5in+0\tabcolsep}}{\advance\leftskip by 5.0pt\relax \advance\rightskip by 5.0pt\relax \textcolor[HTML]{000000}{\fontsize{11}{11}\selectfont{Good\ optimization\ for\ efficient\ analysis,\ integrates\ seamlessly\ with\ ggplot2\ to\ build\ maps}}} & \multicolumn{1}{>{\cellcolor[HTML]{C6FCC2}\raggedright}m{\dimexpr 1.5in+0\tabcolsep}}{\advance\leftskip by 5.0pt\relax \advance\rightskip by 5.0pt\relax \textcolor[HTML]{000000}{\fontsize{11}{11}\selectfont{Doesn't\ work\ the\ best\ on\ Linux\ systems\ (e.g.,\ Mac)}}} \\

\noalign{\global\arrayrulewidth 0pt}\arrayrulecolor[HTML]{000000}





\multicolumn{1}{>{\cellcolor[HTML]{C6FCC2}\raggedright}m{\dimexpr 1.5in+0\tabcolsep}}{\advance\leftskip by 5.0pt\relax \advance\rightskip by 5.0pt\relax \textcolor[HTML]{000000}{\fontsize{11}{11}\selectfont{leaflet}}} & \multicolumn{1}{>{\cellcolor[HTML]{C6FCC2}\raggedright}m{\dimexpr 1.5in+0\tabcolsep}}{\advance\leftskip by 5.0pt\relax \advance\rightskip by 5.0pt\relax \textcolor[HTML]{000000}{\fontsize{11}{11}\selectfont{Interactive\ maps\ for\ HTML}}} & \multicolumn{1}{>{\cellcolor[HTML]{C6FCC2}\raggedright}m{\dimexpr 1.5in+0\tabcolsep}}{\advance\leftskip by 5.0pt\relax \advance\rightskip by 5.0pt\relax \textcolor[HTML]{000000}{\fontsize{11}{11}\selectfont{Build\ attractive\ maps\ with\ toggle-able\ layers}}} & \multicolumn{1}{>{\cellcolor[HTML]{C6FCC2}\raggedright}m{\dimexpr 1.5in+0\tabcolsep}}{\advance\leftskip by 5.0pt\relax \advance\rightskip by 5.0pt\relax \textcolor[HTML]{000000}{\fontsize{11}{11}\selectfont{Limited\ free\ basemaps,\ syntax\ is\ far\ less\ intuitive\ than\ sf}}} \\

\noalign{\global\arrayrulewidth 0pt}\arrayrulecolor[HTML]{000000}





\multicolumn{4}{>{\cellcolor[HTML]{F5F264}\raggedright}m{\dimexpr 6in+6\tabcolsep}}{\advance\leftskip by 5.0pt\relax \advance\rightskip by 5.0pt\relax \textcolor[HTML]{000000}{\fontsize{11}{11}\selectfont{Automation}}} \\

\noalign{\global\arrayrulewidth 0pt}\arrayrulecolor[HTML]{000000}





\multicolumn{1}{>{\cellcolor[HTML]{FAF9AC}\raggedright}m{\dimexpr 1.5in+0\tabcolsep}}{\advance\leftskip by 5.0pt\relax \advance\rightskip by 5.0pt\relax \textcolor[HTML]{000000}{\fontsize{11}{11}\selectfont{futile.logger}}} & \multicolumn{1}{>{\cellcolor[HTML]{FAF9AC}\raggedright}m{\dimexpr 1.5in+0\tabcolsep}}{\advance\leftskip by 5.0pt\relax \advance\rightskip by 5.0pt\relax \textcolor[HTML]{000000}{\fontsize{11}{11}\selectfont{Log\ warnings,\ messages,\ errors\ and\ more}}} & \multicolumn{1}{>{\cellcolor[HTML]{FAF9AC}\raggedright}m{\dimexpr 1.5in+0\tabcolsep}}{\advance\leftskip by 5.0pt\relax \advance\rightskip by 5.0pt\relax \textcolor[HTML]{000000}{\fontsize{11}{11}\selectfont{Great\ for\ logging\ during\ routine\ processes}}} & \multicolumn{1}{>{\cellcolor[HTML]{FAF9AC}\raggedright}m{\dimexpr 1.5in+0\tabcolsep}}{\advance\leftskip by 5.0pt\relax \advance\rightskip by 5.0pt\relax \textcolor[HTML]{000000}{\fontsize{11}{11}\selectfont{None}}} \\

\noalign{\global\arrayrulewidth 0pt}\arrayrulecolor[HTML]{000000}





\multicolumn{1}{>{\cellcolor[HTML]{FAF9AC}\raggedright}m{\dimexpr 1.5in+0\tabcolsep}}{\advance\leftskip by 5.0pt\relax \advance\rightskip by 5.0pt\relax \textcolor[HTML]{000000}{\fontsize{11}{11}\selectfont{tryCatchLog}}} & \multicolumn{1}{>{\cellcolor[HTML]{FAF9AC}\raggedright}m{\dimexpr 1.5in+0\tabcolsep}}{\advance\leftskip by 5.0pt\relax \advance\rightskip by 5.0pt\relax \textcolor[HTML]{000000}{\fontsize{11}{11}\selectfont{Error\ handling}}} & \multicolumn{1}{>{\cellcolor[HTML]{FAF9AC}\raggedright}m{\dimexpr 1.5in+0\tabcolsep}}{\advance\leftskip by 5.0pt\relax \advance\rightskip by 5.0pt\relax \textcolor[HTML]{000000}{\fontsize{11}{11}\selectfont{Can\ be\ used\ to\ prevent\ fatal\ errors\ from\ crashing\ routine\ processes,\ works\ well\ in\ combination\ with\ futile.logger}}} & \multicolumn{1}{>{\cellcolor[HTML]{FAF9AC}\raggedright}m{\dimexpr 1.5in+0\tabcolsep}}{\advance\leftskip by 5.0pt\relax \advance\rightskip by 5.0pt\relax \textcolor[HTML]{000000}{\fontsize{11}{11}\selectfont{Default\ error\ logging\ behaviors\ can\ be\ difficult\ to\ understand}}} \\

\noalign{\global\arrayrulewidth 0pt}\arrayrulecolor[HTML]{000000}





\multicolumn{4}{>{\cellcolor[HTML]{F58771}\raggedright}m{\dimexpr 6in+6\tabcolsep}}{\advance\leftskip by 5.0pt\relax \advance\rightskip by 5.0pt\relax \textcolor[HTML]{000000}{\fontsize{11}{11}\selectfont{HTML}}} \\

\noalign{\global\arrayrulewidth 0pt}\arrayrulecolor[HTML]{000000}





\multicolumn{1}{>{\cellcolor[HTML]{FAC6BB}\raggedright}m{\dimexpr 1.5in+0\tabcolsep}}{\advance\leftskip by 5.0pt\relax \advance\rightskip by 5.0pt\relax \textcolor[HTML]{000000}{\fontsize{11}{11}\selectfont{Bslib}}} & \multicolumn{1}{>{\cellcolor[HTML]{FAC6BB}\raggedright}m{\dimexpr 1.5in+0\tabcolsep}}{\advance\leftskip by 5.0pt\relax \advance\rightskip by 5.0pt\relax \textcolor[HTML]{000000}{\fontsize{11}{11}\selectfont{Modern\ UI\ for\ R\ Shiny\ and\ R\ Markdown\ via\ Bootstrap}}} & \multicolumn{1}{>{\cellcolor[HTML]{FAC6BB}\raggedright}m{\dimexpr 1.5in+0\tabcolsep}}{\advance\leftskip by 5.0pt\relax \advance\rightskip by 5.0pt\relax \textcolor[HTML]{000000}{\fontsize{11}{11}\selectfont{Create\ attractive\ user\ interfaces\ for\ both\ Shiny\ and\ Markdown\ outputs,\ uses\ more\ modern\ version\ of\ Bootstrap}}} & \multicolumn{1}{>{\cellcolor[HTML]{FAC6BB}\raggedright}m{\dimexpr 1.5in+0\tabcolsep}}{\advance\leftskip by 5.0pt\relax \advance\rightskip by 5.0pt\relax \textcolor[HTML]{000000}{\fontsize{11}{11}\selectfont{None}}} \\

\noalign{\global\arrayrulewidth 0pt}\arrayrulecolor[HTML]{000000}





\multicolumn{1}{>{\cellcolor[HTML]{FAC6BB}\raggedright}m{\dimexpr 1.5in+0\tabcolsep}}{\advance\leftskip by 5.0pt\relax \advance\rightskip by 5.0pt\relax \textcolor[HTML]{000000}{\fontsize{11}{11}\selectfont{bookdown}}} & \multicolumn{1}{>{\cellcolor[HTML]{FAC6BB}\raggedright}m{\dimexpr 1.5in+0\tabcolsep}}{\advance\leftskip by 5.0pt\relax \advance\rightskip by 5.0pt\relax \textcolor[HTML]{000000}{\fontsize{11}{11}\selectfont{Write\ HTML,\ PDF,\ and\ ePub\ long-form\ documents}}} & \multicolumn{1}{>{\cellcolor[HTML]{FAC6BB}\raggedright}m{\dimexpr 1.5in+0\tabcolsep}}{\advance\leftskip by 5.0pt\relax \advance\rightskip by 5.0pt\relax \textcolor[HTML]{000000}{\fontsize{11}{11}\selectfont{Relatively\ easy\ formatting\ for\ very\ long\ documents,\ can\ include\ interactive\ features}}} & \multicolumn{1}{>{\cellcolor[HTML]{FAC6BB}\raggedright}m{\dimexpr 1.5in+0\tabcolsep}}{\advance\leftskip by 5.0pt\relax \advance\rightskip by 5.0pt\relax \textcolor[HTML]{000000}{\fontsize{11}{11}\selectfont{Technically\ can\ output\ to\ Word,\ but\ not\ elegantly}}} \\

\noalign{\global\arrayrulewidth 0pt}\arrayrulecolor[HTML]{000000}





\multicolumn{1}{>{\cellcolor[HTML]{FAC6BB}\raggedright}m{\dimexpr 1.5in+0\tabcolsep}}{\advance\leftskip by 5.0pt\relax \advance\rightskip by 5.0pt\relax \textcolor[HTML]{000000}{\fontsize{11}{11}\selectfont{pagedown}}} & \multicolumn{1}{>{\cellcolor[HTML]{FAC6BB}\raggedright}m{\dimexpr 1.5in+0\tabcolsep}}{\advance\leftskip by 5.0pt\relax \advance\rightskip by 5.0pt\relax \textcolor[HTML]{000000}{\fontsize{11}{11}\selectfont{Pagination\ of\ Markdown\ documents}}} & \multicolumn{1}{>{\cellcolor[HTML]{FAC6BB}\raggedright}m{\dimexpr 1.5in+0\tabcolsep}}{\advance\leftskip by 5.0pt\relax \advance\rightskip by 5.0pt\relax \textcolor[HTML]{000000}{\fontsize{11}{11}\selectfont{Makes\ it\ fairly\ easy\ to\ paginate\ Markdown\ and\ convert\ to\ PDF}}} & \multicolumn{1}{>{\cellcolor[HTML]{FAC6BB}\raggedright}m{\dimexpr 1.5in+0\tabcolsep}}{\advance\leftskip by 5.0pt\relax \advance\rightskip by 5.0pt\relax \textcolor[HTML]{000000}{\fontsize{11}{11}\selectfont{Complex\ hTML\ features\ may\ not\ paginate\ properly,\ dependence\ on\ modern\ browser\ access}}} \\

\noalign{\global\arrayrulewidth 0pt}\arrayrulecolor[HTML]{000000}

\hhline{>{\arrayrulecolor[HTML]{666666}\global\arrayrulewidth=1.5pt}->{\arrayrulecolor[HTML]{666666}\global\arrayrulewidth=1.5pt}->{\arrayrulecolor[HTML]{666666}\global\arrayrulewidth=1.5pt}->{\arrayrulecolor[HTML]{666666}\global\arrayrulewidth=1.5pt}-}



\end{longtable}



\arrayrulecolor[HTML]{000000}

\global\setlength{\arrayrulewidth}{\Oldarrayrulewidth}

\global\setlength{\tabcolsep}{\Oldtabcolsep}

\renewcommand*{\arraystretch}{1}

\end{document}
